% !TeX root = ../papers/nc-p3.tex
\section{四次核曲面上の随伴次数と二次錐被覆の母線}
\label{sec:cofactor-eighteen-quartic-and-cone}
\begin{theorem}\label{thm:cofactor-eighteen-normal-quartic}
$d=18$ で四次核曲面上の入力逆写像が余次元一で正則なら矛盾する。
特に正規四次は特異点の種類を問わず核像にならない。
\end{theorem}
\begin{proof}
核写像は双有理で、一般切断 $J$ の算術種数は三または四である。この曲線はGorensteinである。
四なら二次、三次の完全交叉である。実際Castelnuovo等号の六次曲線は二次曲面に入り、
滑らかな二次曲面上では種数から $(3,3)$、二次錐上でも重み六の方程式で三次切断になる。
三なら $J=G$ である。四次曲面 $K$ の特異曲線の各成分は次数三以下である。
一般平面切断は整な平面四次で算術種数三であり、各特異曲線との交点は特異点になるためである。
従って各成分はある二次曲面に含まれる。その上の核グラフ解消内の曲線の入力次数は $2HL=12$ 以下で、
十八次の $C$ を支配しない。基点・孤立特異点上でも同様である。
$K$ の滑らかな開集合の基点外で $J$ は滑らか、それ以外の点では一般入力平面上の階数が三で
$G$ は平面曲線に局所同型なので、この場合も $J$ は局所平面である。

$K$ 上の $J$ のイデアルを $I$ とする。$K$ はGorenstein、$\omega_K=\mathcal O_K$ であり、双対化により
\[
 0\longrightarrow\mathcal O_K\longrightarrow\mathcal Hom_K(I,\mathcal O_K)
 \longrightarrow\omega_J\longrightarrow0 .
\]
右端は余次元一Cohen--Macaulay曲線の $\mathcal Ext^1$ である。
一般入力平面の方程式で割った四入力座標は $\mathcal Hom_K(I,\mathcal O_K)$ の切断となる。
実際有限集合の外で入力射が定義され、逆超平面の零因子は被約な $J$ である。
この開集合で $I$ に掛けて正則関数になる各比は、$\mathcal O_K$ の $S_2$ 性から全ての点でも $I$ を
$\mathcal O_K$ に送る。正規四次曲面なら余次元一での定義は固有な射影空間への付値判定から自動的である。
随伴列でこれらを $\omega_J$ に送ると、その切断比は $J$ の一般点で入力平面の座標を与える
（局所的に $J$ の生成式を掛けてから曲線に制限する）。
正規化上でこの線束の次数は $2p_a(J)-2\le6$ だが、その比は次数八の入力平面八次を双有理に与える。
移動部の次数が八に達し得ず矛盾。
\end{proof}

\begin{theorem}\label{thm:cofactor-eighteen-cone-ruling-degree}
$d=18$ で核像が二次錐なら、核直線束モデル $\mathbb F_2$ の関数体の正規二重被覆 $T$ が
入力の正規化である。錐の負切断と母線を $E,F$ と書くと、入力束は
$g^*(E+3F)$ に数値同値で、一般母線の逆像の入力次数は二である。
分岐因子は $2aE+2bF$ で $a=0,1$, $b=1,2$ に限られる。
\end{theorem}
\begin{proof}
前定理群より次数二被覆 $g:T\to\mathbb F_2$ を考えればよい。正規有限平坦被覆の跡分解を前と同様に用いる。
解消で $U=g^*F$, $L=g^*(E+2F)$、$t=HU$ とする。$UL=2$, $L^2=4$, $HL=6$,
$H g^*E=6-2t\ge0$。$L,U$ の包への直交射影で $8\le6t-t^2$ なので $t=2,3$ である。

$t=3$ なら一般母線逆像は連結であり、入力三次曲線に双有理なので種数零または一である。
種数一は不可能である。実際定理\ref{thm:cofactor-eighteen-no-smooth-quadric}の随伴により
不正則数零の場合は $KH\ge-3$ となり上界 $-4$ に反する。
不正則なら同定理の線織面計算をこの一つの鉛筆に用いる。
正種数底を一族の楕円曲線が支配するので底は楕円曲線であり、初等変換後の射影直線束で
$HF_0=2$、楕円鉛筆の移動像 $U_0$ は $U_0 F_0=1$, $U_0^2=1$ である。
ここで一点だけある鉛筆でも矛盾することを示す。$W=\pi_*\mathcal O(U_0)$ は
楕円底上の階数二、次数一の束である（相対次数一、$\xi^2=\deg W$）。
鉛筆の二切断は一般独立であり、比例なら共通の水平零因子を持つ。
従って $0\to\mathcal O^2\to W\to k_p\to0$。二つの定数基底を変えれば非零な拡大類は
$(1,0)$ となり $W=\mathcal O(p)\oplus\mathcal O$ である。
負切断 $C_0$ の自乗は $-1$、$H^2=8$, $HF_0=2$ から $HC_0=1$。
楕円曲線の次数一の線束は基点なしでないので、入力系の生成性に反する。

従って $t=3$ なら一般母線逆像は有理で、分岐因子は $2E+2bF$ である。
$T$ は有理曲面、二重被覆列とRiemann--Rochから $\chi(\mathcal O_T)=1$ なので特異点は有理二重点。
最小解消の標準類は $K=-L+(b-2)U$ である。このモデルでの入力移動部にさらに基点があれば
解消後の $KH$ はその重複度和だけ増える。従って $3b-12\le-4$。
分岐の被約性から $b\ge1$（$0$ 以下では $E$ が二重になる）なので $b=1,2$。

$b=2$ では $K=-L$ である。最小解消上の入力束の自乗は八か九で、随伴の偶奇から八、基点なし。
定理\ref{thm:cofactor-eighteen-no-smooth-quadric}の次数四del Pezzoの議論がそのまま使える。
すなわち $H=L+P$ で $P$ は自由な有理鉛筆、跡切断の零と極は $L$ 次数零の曲線に限られる。
頂点上の曲線は全て一般入力切断を避ける
（$Hg^*E=0$。$E$ 上の点の上の解消例外も全引戻しの成分である）。
それ以外では核曲面は滑らかで、特異被覆点 $y$ を通る核超平面の固定サイクル $Z_y$ は
$(L-Z_y)^2\ge2$ より $HZ_y\le1$。従って前記定理と同じ繊維ごとの評価から
一般平面でEuler数五と四が等しくなり矛盾する。

$b=1$ では分岐は $E$ とそれに交わらない $E+2F$ の切断の和である
（$E$ は固定成分で、残余の水平切断に垂直成分を加えると再び $E$ を含む）。
二切断で分岐する母線ごとの二重被覆は $\mathbb F_1$ である。
実際二切断を零と無限に取った直線束の重み二の遷移関数の平方根を用いればよい。
負切断を $s$、母線を $f$ として $L=2s+2f=2h$、$h=s+f$ は平面への縮約。
交点 $Hf=3$, $HL=6$ を用いると、解消前の系は $3h$ で一単純基点がある。
$s$ は固定でないので基点は $s$ の外。入力正規化は、平面をこの基点で吹いた曲面で
$H=3h-e$ は豊富、入力は有限である。核写像はこの平面上の四つの二次式による
次数二の射と合成される。縮約した $s$ の像以外では、この平面から錐への射は
各点で微分階数一以上である（滑らかな母線方向二重写像と底の座標）。
一般平面切断は $e$ に横断的に一度会い、基点での接方向は $e$ 上を動くので
核写像の微分もその切断方向には非零である。
$e$ の一般入力像は階数三で、同じ核点を持つ他の平面点は高々一つ、その固定入力像を一般平面は避ける。
従って $e$ の一般点の入力近傍で有限正規化は同型である。具体的に $(f,g)$ はそこで
接方向を $f|_e$ と法方向を核座標で分離し、他の逆像はなく、階数三では核は入力から正則に復元される。
しかし跡の有効零因子 $2L-K-2H=h+e$ は $e$ との交点が $-1$ で $e$ を含む必要がある。
同型点では跡合成は単元なので矛盾。

残る $t=2$ はHodgeの等号で $H\equiv L+U$ である。$E+3F$ は豊富だから入力正規化は $T$ で射は有限。
一般逆像が連結なら入力円錐曲線に双有理で種数零、非連結なら一般母線に分岐がない。
従って $a=1,0$ の二択である。有限跡の有効因子は数値的に
$2L-K_T-2H\equiv g^*((2-a)E+(2-b)F)$ である。その押出しの有効性から $b\le2$。
$a=1$ なら上記の被約性、$a=0$ なら被覆の整性より $b\ge1$。
\end{proof}
\path{computations/verify_cofactor_cone_ruling_degree.py} はHodgeの等号、被覆の標準類と残る分岐次数を検算する。
